\section{Vnitřní reprezentace}

Překladač využívá pro reprezentaci zpracovávaného programu dvě vnitřní struktury:
\begin{itemize}
    \item \textbf{Abstraktní syntaktický strom (\texttt{ast}):} Stromová reprezentace programu tvořená uzly (\texttt{Node}), kterou vytváří parser na základě syntaktické, precedenční a sémantické analýzy. Zachytává hierarchii funkcí, příkazů, volání (včetně výrazů v argumentech díky rozšíření \texttt{FUNEXP}) a výrazů.
    \item \textbf{Seznam instrukcí tříadresného kódu (\texttt{ilist}):} Lineární seznam instrukcí vnitřního kódu, do kterého se postupně ukládají vygenerované operace a který slouží jako hlavní datový vstup pro generátor cílového kódu.
\end{itemize}

\subsection{Generování vnitřního kódu}
Převod abstraktního syntaktického stromu na seznam instrukcí tříadresného kódu zajišťuje modul \texttt{insidecodegene}. Tento modul je implementován odděleně od samotných datových struktur \texttt{ast} a \texttt{ilist}.
   Při průchodu uzly \texttt{AST} generátor plní strukturu \texttt{ilist} instrukcemi tříadresného kódu a zároveň řeší:
\begin{itemize}
    \item \textbf{Typy výrazů:} určení datových typů při překladu výrazů a mezivýsledků,
    \item \textbf{Návratové hodnoty:} sjednocení práce s návratovými hodnotami,
    \item \textbf{Vestavěné funkce:} vygenerování instrukcí pro funkce typu \texttt{subStr} či \texttt{ord} a související stavy.
\end{itemize}
    

Výsledkem generování je list instrukcí 3 adresného kódu, který je pak dále předán modulu pro generování cílového kódu. 